\section{Related Work}\label{sec:related}

\myparagraph{Video Generation.} The field of video generation has seen rapid advancements, largely building upon the success of diffusion models in visual synthesis~\cite{ho2020denoising, song2020denoising}. Current approaches can be broadly categorized based on their underlying architecture.
The first category consists of methods based on pre-trained text-to-image U-Net models~\cite{esser2024scaling, chen2024pixart}. These approaches typically insert temporal modules, such as attention or convolution layers, into the frozen U-Net backbone and then fine-tune the model on video data~\cite{guo2023animatediff, wang2023modelscope}. This allows them to leverage powerful image priors, but their capabilities can be constrained by the original image-centric architecture.
The second category is represented by models that adopt a Diffusion Transformer (DiT) architecture~\cite{brooks2024video, yang2024cogvideox, zheng2024open, kong2024hunyuanvideo, ma2025step, polyak2024movie,chen2024gentron,menapace2024snap}. These methods treat video as a sequence of spatiotemporal patches, processing them in a unified manner with a Transformer. This approach has demonstrated superior scalability and flexibility, enabling the generation of high-resolution videos with variable durations and aspect ratios, especially when trained on massive datasets.
The third category explores the integration of Large Language Models (LLMs) with diffusion models. In the image generation domain, works such as~\cite{yu2023scaling, koh2023generating,pan2025transfer,wu2025qwen,wu2024next,shi2024lmfusion}have shown that LLMs can enhance compositional understanding and planning. For video generation, this area is still in a nascent stage but holds significant promise for improving the logical coherence and narrative structure of generated content.

\myparagraph{Video Avatar Model.} Video avatars aims to create realistic human videos from various driving signals, with methods typically categorized by their driving source and synthesis pipeline. One category is pose-driven animation, which generates video based on an explicit, externally provided motion sequence, such as skeletal poses~\cite{zhang2024mimicmotion, shao2024human4dit,chang2023magicpose,tu2024motionfollower,wang2024unianimate,karras2023dreampose,xu2024magicanimate}. As the motion generation step is largely bypassed in these methods, their focus shifts to achieving high-fidelity rendering. More central to this paper is the second key category: audio-driven animation. This task requires a model to first generate plausible human motion from an audio signal and subsequently synthesize the final video. Consequently, the model must handle the dual challenges of both motion generation and rendering. Within this category, a common approach is a two-stage pipeline, where a motion generation model first translates audio into an intermediate representation, such as 2D/3D keypoints or 3D mesh sequences~\cite{zhuang2024vlogger, meng2025echomimicv2,deng2025stereo,wei2024aniportrait} or discrete motion tokens~\cite{hogue2024diffted, tian2025emo2}. A subsequent rendering model then synthesizes the video conditioned on this motion. More recently, end-to-end approaches have emerged~\cite{lin2025cyberhost,lin2025omnihuman,wang2025interacthuman,liang2025alignhuman}, which directly generate video from audio in a single step, aiming for improved audio-motion synchronization.
Despite their architectural differences, these audio-driven methods predominantly treat motion generation as a direct mapping process, a fast, reactive function from audio features to motion output. This process does not explicitly model the high-level cognitive phase where humans' planning and reasoning ultimately determine the resulting physical actions. We believe that incorporating such a reasoning phase is essential for generating more plausible and intelligent human behaviors, and our work aims to provide a foundational exploration in this direction.


\myparagraph{Large Language Models on Cognitive Simulation.} Large Language Models (LLMs) have reshaped machine reasoning with a clear upward trajectory in cognitive capability. Foundational models like GPT-3~\cite{brown2020language} and ChatGPT~\cite{ouyang2022training, achiam2023gpt} laid the groundwork, while recent iterations such as GPT-4o~\cite{hurst2024gpt}, OpenAI o-series~\cite{jaech2024openai} and DeepSeek-R1~\cite{guo2025deepseek} enhance contextual and multi-modal reasoning, with prompting being key to unlocking this potential—Chain-of-Thought (CoT)~\cite{wei2022chain} revolutionized problem-solving by decomposing tasks into logical steps. Derivatives like Self-Consistency~\cite{wangself} using aggregating reasoning paths and Tree-of-Thought (ToT)~\cite{yao2023tree} with branching exploration further improved accuracy in arithmetic, logic, and commonsense tasks, confirming LLMs can simulate human-like deductive/inductive thinking via structured prompting. 
This reasoning capability has further empowered autonomous agents as their core engine. Classic practical frameworks like Toolformer~\cite{schick2023toolformer}, MetaGPT~\cite{hong2023metagpt}, and AutoGPT~\cite{zhang2023autogpt} with autonomous goal refinement enable agents to transcend pre-defined rules, exhibiting intent understanding and error correction. Voyager~\cite{wang2023voyager} leverages LLMs for open-world strategies and adaptive responses, outperforming rule-based agents. Furthermore, the renowned Generative Agents~\cite{park2023generative} enable LLMs to simulate daily human behaviors (e.g., social interaction, scheduling), marking a pivotal shift in human-like agent simulation. Recent advancements like OpenAI's Deep Research~\cite{openai2025deepresearch} and Monica's Manus~\cite{monica2025manus} further showcase LLM-driven proficiency in complex task planning and cross-domain execution.
Beyond agent systems, LLM reasoning also steers generative models across domains to solve the "controllable generation" challenge. For image editing, InstructPix2Pix~\cite{brooks2023instructpix2pix} relies on LLM-generated instructions to refine image modifications. MetaQueries~\cite{pan2025transfer} uses learnable queries to prompt vision-language models to guide image generation with enhanced semantic alignment. Besides, recent LLM-driven video generation agents~\cite{hu2024storyagent, yuan2024mora, li2024anim, wu2025automated} enable more controllable long video synthesis through collaborative agentic workflows. These works collectively position LLMs as universal planners that bridge high-level user intent and the execution logic of low-level generative models. Yet notably, the integration of LLM-driven reasoning and planning into more fine-grained intelligent human/avatar behavior generation remains an underexplored area.
